% GENERATED by scripts/build-lecture-slides.py; do not edit. Portable profile: XeTeX, theme and fonts from the core.
\makeatletter\IfFileExists{core/beamer/beamerthememx.sty}{\def\mxcorepath{core/}\def\input@path{{core/beamer/}{core/}}}{\def\mxcorepath{../core/}\def\input@path{{../core/beamer/}{../core/}}}\makeatother
\documentclass[aspectratio=169,10pt,t]{beamer}
\usetheme{mx}
\usepackage{graphicx,booktabs}
\usepackage[french]{babel}
\setlength{\parskip}{.45em}
\hypersetup{pdfauthor={},pdfsubject={Lecture slides},pdfkeywords={},pdftitle={Trigonométrie}}
\providecommand{\N}{\mathbb{N}}
\providecommand{\Z}{\mathbb{Z}}
\providecommand{\D}{\mathbb{D}}
\providecommand{\Q}{\mathbb{Q}}
\providecommand{\R}{\mathbb{R}}
\providecommand{\C}{\mathbb{C}}
\providecommand{\vect}[1]{\overrightarrow{#1}}
\providecommand{\vi}{\vec{\imath}}
\providecommand{\vj}{\vec{\jmath}}
\providecommand{\vk}{\vec{k}}
\providecommand{\norm}[1]{\left\lVert #1\right\rVert}
\providecommand{\vc}[2]{\begin{pmatrix}#1\\#2\end{pmatrix}}
\providecommand{\deter}[2]{\det\left(#1,#2\right)}
\providecommand{\abs}[1]{\left\lvert #1\right\rvert}
\providecommand{\dd}{\,\mathrm{d}}
\providecommand{\e}{\mathrm{e}}
\providecommand{\Card}{\operatorname{Card}}
\providecommand{\E}{\operatorname{E}}
\providecommand{\V}{\operatorname{V}}
\providecommand{\degre}{^{\circ}}
\providecommand{\un}[1]{\,\mathrm{#1}}
\providecommand{\sci}[2]{#1\times10^{#2}}
\providecommand{\pgcd}{\operatorname{PGCD}}
\providecommand{\ppcm}{\operatorname{PPCM}}
\title{Trigonométrie}
\subtitle{Cours et exemples}
\date{}
\author{}
\begin{document}
\mxtitleframe
\begin{frame}{Objectifs}
Convertir les unités d’angle, transformer des expressions et résoudre des équations et inéquations trigonométriques.
\end{frame}
\begin{frame}{Mesures et cercle orienté}
Un tour correspond à \(360^{\circ}=2\pi\) radians \(=400\) grades. Les mesures d’un angle orienté diffèrent de \(2k\pi\) (\(k\in\mathbb{Z}\)). Sa mesure principale est l’unique représentant dans \(]-\pi;\pi]\).
\end{frame}
\begin{frame}{Mesures et cercle orienté}
\(7\pi/3\) a mesure principale \(\pi/3\). \(270^{\circ}=3\pi/2=300\) grades; la mesure principale en radians est \(-\pi/2\).
\end{frame}
\begin{frame}{Formules à utiliser et à justifier}
\(\cos(-x)=\cos x\), \(\sin(-x)=-\sin x\);
\(\cos(\pi-x)=-\cos x\), \(\sin(\pi-x)=\sin x\);
\(\cos(x+\pi)=-\cos x\), \(\sin(x+\pi)=-\sin x\).
\[
\begin{aligned}
\cos(a+b)&=\cos a\cos b-\sin a\sin b,\\
\sin(a+b)&=\sin a\cos b+\cos a\sin b,\\
\cos(a-b)&=\cos a\cos b+\sin a\sin b,\\
\sin(a-b)&=\sin a\cos b-\cos a\sin b.
\end{aligned}
\]
\end{frame}
\begin{frame}{Formules à utiliser et à justifier}
Les symétries du cercle donnent les angles associés. Pour la formule du cosinus, prendre les vecteurs unitaires \((\cos a;\sin a)\) et \((\cos(-b);\sin(-b))\). Leur produit scalaire vaut à la fois \(\cos(a+b)\) et \(\cos a\cos b-\sin a\sin b\). La formule du sinus suit de \(\sin(a+b)=\cos(\pi/2-a-b)\), en utilisant les coordonnées des angles complémentaires. Les formules de différence suivent en remplaçant \(b\) par \(-b\).
\end{frame}
\begin{frame}{Formules à utiliser et à justifier}
\(\sin2x=2\sin x\cos x\);
\(\cos2x=\cos^2x-\sin^2x=2\cos^2x-1=1-2\sin^2x\).
\[
\begin{aligned}
\cos a+\cos b&=2\cos\tfrac{a+b}{2}\cos\tfrac{a-b}{2},\\
\sin a+\sin b&=2\sin\tfrac{a+b}{2}\cos\tfrac{a-b}{2},\\
2\cos a\cos b&=\cos(a+b)+\cos(a-b),\\
2\sin a\sin b&=\cos(a-b)-\cos(a+b).
\end{aligned}
\]
\end{frame}
\begin{frame}{Formules à utiliser et à justifier}
La duplication prend \(a=b=x\). Pour les sommes, poser \(u=(a+b)/2,v=(a-b)/2\), puis additionner les formules en \(u+v\) et \(u-v\). Les produits résultent de l’addition ou de la soustraction des formules d’addition.
\end{frame}
\begin{frame}{Équations et inéquations}
Pour \(k\in\mathbb{Z}\),
\[
\cos x=\cos a\Leftrightarrow x=2k\pi\pm a,
\]
\[
\sin x=\sin a\Leftrightarrow x=a+2k\pi\ \text{ou}\ x=\pi-a+2k\pi.
\]
Pour les tangentes définies, \(\tan x=\tan a\Leftrightarrow x=a+k\pi\).
\end{frame}
\begin{frame}{Équations et inéquations}
Une abscisse sur le cercle correspond aux deux points symétriques par rapport à l’axe horizontal; une ordonnée correspond aux points symétriques par rapport à l’axe vertical. Ajouter les tours complets donne toutes les mesures. Si les deux cosinus sont non nuls, l’égalité des tangentes équivaut à \(\sin x\cos a-\cos x\sin a=0\), soit \(\sin(x-a)=0\); donc \(x-a=k\pi\).
\end{frame}
\begin{frame}{Équations et inéquations}
Dessiner les intersections de la droite de niveau avec le cercle; identifier les arcs qui satisfont l’inégalité; traduire les arcs en intervalles; conserver les bornes selon le signe strict ou large. Pour \(\tan\), exclure les zéros de \(\cos\).
\end{frame}
\begin{frame}{Équations et inéquations}
Sur \([0;2\pi[\), \(\cos x\ge1/2\) donne \([0;\pi/3]\cup[5\pi/3;2\pi[\). Dans \(\mathbb{R}\), la solution est la réunion des intervalles \([-\pi/3+2k\pi;\pi/3+2k\pi]\).
\end{frame}
\begin{frame}{Équations et inéquations}
Poser provisoirement \(X=\sin x\) ou \(X=\cos x\), résoudre l’équation algébrique puis garder seulement les racines dans \([-1;1]\). Pour chaque valeur admissible, écrire toutes les familles trigonométriques. Enfin restreindre à l’intervalle demandé et retirer les doublons. Poser \(X\) ne dispense pas de revenir à \(x\).
\end{frame}
\begin{frame}{Équations et inéquations}
Pour \(\cos(2x)=1/2\) et \(x\in[0;\pi]\), l’angle \(2x\) appartient à \([0;2\pi]\). Ses valeurs possibles sont \(\pi/3\) et \(5\pi/3\), donc \(x=\pi/6\) ou \(5\pi/6\). Dans \(\mathbb{R}\), les familles sont \(x=k\pi\pm\pi/6\); l’intervalle est imposé après avoir divisé toutes les mesures, y compris la période, par \(2\).
\end{frame}
\begin{frame}{Équations et inéquations}
Une consigne avec « et » demande l’intersection des ensembles solutions; « ou » demande leur réunion. Sur un même cercle, marquer les arcs, puis vérifier les bornes une par une. Une inégalité stricte exclut ses points d’égalité, même s’ils sont admis dans l’autre condition.
\end{frame}
\begin{frame}{Équations et inéquations}
Écrire \(\sin x=\sin a\) ne permet pas de conclure \(x=a\) : il y a une seconde famille et des tours complets. De même, la calculatrice fournit souvent une valeur principale, pas l’ensemble des solutions d’une équation.
\end{frame}
\end{document}
